\documentclass[11pt]{article}
\usepackage[margin=27mm]{geometry}
\usepackage{amsmath,amssymb,amsthm,booktabs,longtable,hyperref}
\setlength{\emergencystretch}{2em}
\newtheorem{theorem}{Theorem}\newtheorem{lemma}[theorem]{Lemma}
\newcommand{\PSL}{\operatorname{PSL}}\newcommand{\PGL}{\operatorname{PGL}}
\newcommand{\Aut}{\operatorname{Aut}}\newcommand{\F}{\mathbb F}
\newcommand{\N}{\operatorname{N}}
\title{Rank two projective geometries: classification and corrections}
\author{Human Compression Project}\date{}
\begin{document}\maketitle
\begin{abstract}
We give a complete treatment of the mathematical claims in the fixed
published Part II article and its retained appendix. The published
classification omits affine configurations over subfields. Its convention
identifies the two parabolic subgroups without ordering them, which also
changes two families of conjugacy counts. Four stated full graph groups
require correction. The surviving configurations, their class counts,
their local arc lengths and the remaining full group assertions follow
from the proofs collected here. We distinguish these corrections from
empty parameter cases, preliminary slips, classical premises and the
contents missing from the retained appendix edition.
\end{abstract}

\section{The two originals and the meaning of completion}
The first original, designated L063, is Buekenhout, De Saedeleer and
Leemans, \emph{On the rank two geometries of the groups $\PSL(2,q)$:
part II}, Ars Mathematica Contemporanea 6 (2013), 365--388. It has
24 pages. The second, L064, is a retained 21-page appendix edition
beginning with the proof of Lemma 12 and ending with Proposition 30.
They are two separately registered originals. Shared arguments are
included once in this collection, while the claim ledgers identify
each original and its actual wording separately.

A complete treatment of a false assertion is its explicit disproof,
not an attribution of the corrected assertion to the source. The omitted
affine family was already established in the earlier accepted component;
it is reused here. The present completion adds the proofs needed for the
surviving assertions. It does not record another principal refutation of
either original. In particular, the source's four impossible square-field
rows are empty conditional cases, rather than four numerical counterexamples.

Our scope is the mathematical classification and its stated consequences.
Historical computer timings and statements that particular examples were
checked in MAGMA are attributed to their authors; those historical runs
are not reproduced. The classification does not depend on them. The
authors' earlier constructions, the Biggs--Hoare and Wong graphs, and the
classical generalized quadrangle retain their original attribution.

\section{A common proof structure}
Put $G=\PSL(2,q)$, and write $\{M,L\}$ for two distinct proper
parabolic subgroups, with $B=M\cap L$. We use the unordered convention
in Section 2 of L063. Conjugacy means an orbit under $G$; isomorphism
in that source means an orbit under $\Aut(G)$. This latter convention
concerns the distinguished action and does not assert that every abstract
graph isomorphism is induced by $\Aut(G)$.

The coset construction gives one object for each coset of $M$ or $L$,
with incidence defined by nonempty intersection. Every incident pair
is translated to the standard pair, whose stabilizer is $B$. Its local
actions are $M$ on $M/B$ and $L$ on $L/B$. Thus the required rank two
conditions amount to
\[
 M\text{ or }L\text{ maximal in }G,
 \qquad |B\backslash M/B|=|B\backslash L/B|=2.
\]
The local actions are primitive and have degree at least two, proving
firmness. The other parabolic is not contained in the maximal one, so
the two generate $G$; this proves connectedness. For fixed $M$, ordered
configurations are normalizer orbits on $L$. If the parabolic types agree,
one must also identify the orbits obtained by exchanging them. This is
the accepted common rank two argument from the Part I treatment, restated
to fix the convention used throughout this collection.

The next step is to classify the possible two-transitive actions. For
$d=[H:B]$, a point stabilizer is transitive on the other $d-1$ points,
so
\[
 d-1\mid |B|,\qquad |H|\le |B|(|B|+1).
\]
The separately proved common action theorem combines this bound with
the classical subgroup types, and treats the normal $\PSL$ subgroup of
$\PGL$ explicitly. It supplies the projective, affine and small
exceptional possibilities used in every subsequent application. We do
not assume the source's representation proposition, whose cited proof
was listed as being in preparation.

\begin{center}
\begin{tabular}{p{.35\textwidth}p{.56\textwidth}}
\toprule
Proof component & Its role in the classification\\\midrule
Two-transitive coset actions in dimension two & Complete candidate lists
for the local parabolics and intersections.\\
Subfield planes and unordered pairs of parabolics & Affine planes, their
conjugacy and Frobenius orbits, and the odd-characteristic unordered
count.\\
Counting configurations above a cyclic subgroup of order three & All seven Table 2 rows,
including the corrected equal-type inner count.\\
Exceptional parabolics from their common normalizers & Tables 1 and 3
exceptional counts and the empty square-field cases.\\
Subfield parabolics and affine overgroups & Complete even, odd and
quadratic subfield alternatives, including the small fields.\\
Local arc transitivity from small stabilizers & Every local arc column
in Tables 1--6, including the nine Table 9 crossreferences.\\
The three graph group components & Full groups for the $S_4,S_4$
families, the $S_4,D_{16}$ subdivisions, and the $q=9$ quadrangle.\\
\bottomrule
\end{tabular}
\end{center}

The subgroup classification of Dickson, Moore and Suzuki is an explicit
classical premise, including the stated maximal subgroups, subfield
embeddings, normalizers and fusion under $\PGL$. The order of $\PSL$
and its action on a projective line are classical as well. The full graph
group proofs additionally use Tutte's bound of 48 for a vertex stabilizer
in a finite connected cubic arc-transitive graph, under its actual
hypotheses. None of these premises is a retained version of the target
Part II geometry classification. The accepted Part I treatment is used
only where a maximal parabolic belongs to its Borel or dihedral scope.

\section{The corrected classification and numerical columns}
The complete action list separates the maximal parabolics of the Part II
problem into $A_4,S_4,A_5$, proper $\PSL$ subfield groups and quadratic
$\PGL$ subfield groups. For each allowed intersection, the common
normalizer proofs count actual configurations, rather than inferring
their number merely from subgroup conjugacy. The exceptional and small
field arguments supply all finite branches. The subfield proof shows
that there are no further branches.

For clarity, the changes to the printed numerical columns are the
following. All other nonempty rows retain their source class counts.
\begin{enumerate}
\item In even characteristic, let $K=\F_s\subset F=\F_{s^m}$ with
$s\ge4$ and $m$ prime. Beside $M=\PSL(2,K)$, take
\[
 B=U_K:K^*,\qquad L=U_V:K^*,
 \qquad K\subset V\subset F,\quad\dim_KV=2.
\]
Every such $V$ gives the required configuration. There are
$(s^{m-1}-1)/(s-1)$ conjugacy classes, and the automorphism classes
are exactly the Frobenius orbits on these subspaces. At $m=2$ this
is the single printed Table 4 row. The assertion that $m=2$ is
necessary is false. At $s=4$ the accepted scalar formulas recover
the Table 1 counts as well.
\item In Table 2 rows 6--7, write $m=(p-\epsilon)/6$, where
$3\mid(p-\epsilon)/2$. The printed $2m-1$ counts ordered
configurations. For unordered configurations the correct inner count is
\[
 m+\lfloor m/2\rfloor.
\]
The automorphism count remains $m$. The proof identifies the
overgroups with a regular dihedral normalizer quotient and then
takes inversion orbits.
\item In Table 5, for $q=3^m$ with $m$ an odd prime, the inner count
is $(3^{m-1}-1)/2$, rather than $3^{m-1}-1$. At $q=27$ it is
four rather than eight. The printed automorphism counts are retained:
two at $m=3$, and $(3^{m-1}-1)/(2m)$ otherwise.
\item In the square-field part of Table 1, $q=p^2\equiv-1\pmod5$
with $p$ odd implies $q\equiv9\pmod{40}$. Rows 12--14 and 18
therefore have no admissible parameters. Rows 15--17 remain; the
alternative $q\equiv-1\pmod{40}$ in rows 16--17 is also empty.
\end{enumerate}

The complete source tables have $18+7+9+5+2+3=44$ rows. This is a
count of printed rows, not of distinct families or new theorems. Table 4
row 5 repeats Table 1 row 1; the $s=4$ subfield case overlaps the
$A_5$ discussion; four printed square-field rows are empty. The ledgers
record these overlaps without adding source credit. The nine entries of
Table 9 point back to those same rows.

The original local arc lengths are proved for each actual configuration.
The added even-characteristic affine configurations have local arc length
three by the same local calculation. The $S_4,S_4,D_8$ configuration
has length four for the specified simple group, and the
$S_4,D_{16},D_8$ configuration has length seven. These statements
about $G$ are kept separate from the full graph groups.

\section{Full graph groups}
For a prime $p\equiv\pm1\pmod8$, choose $\epsilon\in\{1,-1\}$
with $p\equiv\epsilon\pmod8$, and put $d=(p-\epsilon)/8$.
The full group of a graph with parabolics $S_4,S_4$ and intersection
$D_8$ is
\[
 \Aut X\cong
 \begin{cases}
 \PSL(2,p)\times C_2,&d\text{ even},\\
 \PGL(2,p),&d\text{ odd}.
 \end{cases}
\]
An explicit torus-normalizer calculation constructs the required type
interchange. Tutte's bound then gives index at most four over $G$.
Simplicity forces $G$ to be normal, and the centralizer and the action
on the two $S_4$ classes determine which of the displayed groups occurs.
Thus the printed $\PGL$ identifications at $p=17,31,47,79$ are
incorrect; the identifications at $p=23,41,71,73,89$ are recovered.
In all these prime cases the full graph action still has arc length four.

For $q=9$, the graph has the 15 duads and 15 synthemes on six letters
as its two types. Its symplectic realization proves that it is the
generalized quadrangle of order $(2,2)$. A displayed permutation
interchanges types and is checked directly on all 45 edges. Together
with the faithful $S_6$ action this attains the Tutte bound of 1440.
The full group is therefore $\operatorname{P\Gamma L}(2,9)$, with full arc length
five, as the source states. The construction and the named graph are
classical; the literal duality certificate is a verification of the proof.

For $q=17,31,79,97$, suppressing the degree-two vertices in the
$S_4,D_{16}$ graph gives a simple cubic graph. The suppression proof
excludes parallel edges and uses the actual transitive action. Tutte's
bound gives index at most two, while the outer automorphism exchanges
the two $S_4$ classes and cannot preserve the required action.
The full group is $\PSL(2,q)$, recovering all four stated groups.

\section{Preliminary formulas and literal corrections}
The elementary order obstruction above proves L063 Lemma 4.4.
For Lemma 4.5, an abelian transitive permutation group has the same
stabilizer at every point. Faithfulness makes that stabilizer trivial,
so the action is regular. If it is two-transitive, its trivial point
stabilizer can be transitive on the other points only when there is
one other point. Thus the degree is two. The word ``permutation''
here includes faithfulness; a group with a nonfaithful abelian action
would first be replaced by its image.

The elementary abelian orbit formulas in Lemma 4.1 need their field
parameter to be specified. This also explains why the source correctly
declines to give a number of conjugacy classes from those formulas alone.
\begin{lemma}\label{multiplier}
Let $0\ne V\subset\F_q$ be an additive subgroup of order $p^m$,
and let
\[
 K=\{a\in\F_q:aV\subset V\}=\F_{p^k}.
\]
Put $A_F=\F_q^*$ in characteristic two and $A_F=(\F_q^*)^2$
in odd characteristic. In $G=\PSL(2,q)$,
\[
 \N_G(U_V)=U_{\F_q}:(K^*\cap A_F).
\]
If $1\ne D\le K^*\cap A_F$, then
\[
 \N_G(U_V:D)=U_V:(K^*\cap A_F).
\]
Consequently the orbit lengths are respectively
\[
 \frac{q^2-1}{e(p^k-1)},\qquad
 \frac{(q^2-1)p^{n-m}}{e(p^k-1)},
 \quad q=p^n,
\]
where $e=2$ if $p$ is odd and $n/k$ is even, and $e=1$ otherwise.
\end{lemma}
\begin{proof}
The set $K$ is closed under addition and multiplication and contains
one. Every nonzero element sends the finite set $V$ bijectively to
itself, so its inverse also lies in $K$. Thus $K$ is a field, and
$k$ divides both $n$ and $m$. A normalizer must fix the unique
fixed point of $U_V$, hence has affine form $x\mapsto ax+b$.
It normalizes $U_V$ exactly when $a\in K^*\cap A_F$; the
translation $b$ is unrestricted. For the second group $U_V$ is
characteristic as its unique Sylow subgroup for the characteristic.
Conjugating any scalar $c\ne1$ in $D$ gives translation
$(1-c)b$, which lies in $V$ exactly when $b\in V$, since
$(1-c)^{-1}\in K$. The displayed normalizers follow. Divide
$|G|=q(q^2-1)/\gcd(2,q-1)$ by their orders and use the square
criterion for $K^*$ in $\F_q$ to obtain the formulas.
\end{proof}
Thus the source's $k$ is the degree of this multiplier field, not an
arbitrary common divisor. Different $V$ can have different such fields;
the formulas determine each orbit length, not the number of orbits.

Three additional preliminary corrections are explicit.
\begin{enumerate}
\item The odd-characteristic single-class clauses of Lemma 4.1(5)
must require $d>2$ when referring to $D_{2d}$. At $d=2$ the
four-groups belong to item (6). For example, at $q=5$ the former
formula would give 15 four-groups, whereas $A_5$ has five: each
is the subgroup of double transpositions fixing one of the five
letters. This verifies the correction without an ambient search.
\item The $\PGL(2,s)$ index in Table 8 is
\[
 [\PSL(2,s^2):\PGL(2,s)]
 =\frac{q(q^2-1)}{2s(s^2-1)},\qquad q=s^2.
\]
The printed denominator omits the factor two. This follows directly
from the two group orders.
\item The printed exhaustive list in Proposition 4.6 omits the
degree-two and degree-four actions of a maximal $S_4$ at prime $q$.
For instance $S_4<\PSL(2,7)$ has its action on $S_4/A_4$ of
degree two and its faithful natural action on $S_4/S_3$ of degree
four. Its standalone clause lists degree three only; the separate
$\PGL$ clauses impose $q=s^2$. The common action proof includes
all three $S_4$ actions with their correct kernels. Assertions about
representations inside one fixed subgroup are distinguished from
the two ambient embedding classes.
\end{enumerate}

L064's proof of Lemma 24 also incorrectly excludes every
$\PGL(2,3^w)$ from $\PSL(2,3^{2m})$ for odd prime $m$.
Both $w=1$ and $w=m$ give subfield embeddings. The required
two-transitive overgroup exclusion is nevertheless true and is proved
by the common action theorem; it does not use that auxiliary assertion.
In L063's proof of Lemma 5.3, Steps 4--5 confuse the labels $A_4$
and $E_{16}:3$. The total number of the $A_4$ subgroups in the
stated characteristic-two family is $q(q^2-1)/12$, rather than
$q(q-1)/12$: the additional factor $q+1$ counts their unique
fixed projective points. The accepted coordinate proof gives the
correct intermediate counts while recovering the lemma's final formula.
In three displayed transporter formulas of L064 Proposition 30,
on physical pages 19, 20 and 21, the denominator must be the actual
ambient order $|\PSL(2,p^{2n})|$, not
$|\PSL(2,2^{2n})|$. The literal normalizer proofs give the stated
final values without that transcription error.
In Section 6 of L063, the second path in the $A_5,D_{30}$ argument
starts at type one, and the $S_4,D_{24}$ argument must say $D_{24}$,
not $D_{18}$, when acting on cosets of $D_8$. The full local proofs
in the collection fix both slips.

\section{Separate attribution and source accounting}
L063's Theorem 1.1 and its concluding completeness assertion are false
as printed because of the omitted affine family. The corrected proof
retains the valid configurations and all valid class counts, then
replaces the incomplete subfield argument and the unordered counts.
The local arc columns are recovered, and the full graph groups have
the stated four corrections. Preliminary borrowed results are identified
as premises with the exact qualifications above; we do not claim to
have rederived the classical subgroup classification.

L064 contains proofs for Lemma 12, Proposition 16, Lemmas 17--19,
Proposition 20, Lemma 21, Proposition 22, Lemmas 23--25,
Proposition 26, Lemmas 27--29 and Proposition 30. Its Lemma 21
lists the affine type without the published lemma's restriction, but
Proposition 22 still restricts the configuration to extension degree
two. The respective dispositions are therefore different. Its
Propositions 16 and 26 share the unordered-count corrections, and
Lemma 24 has the auxiliary correction just discussed. Its other
surviving assertions follow from the same explicit arguments as L063.

The fixed appendix does not contain the advertised entire $A_5$ count
proof, the standalone published $q=25$/$S_4$ lemma, or Section 6's
claimed arc proofs on appendix pages 19--20. Those pages actually
concern Proposition 30. We supply the required proofs for L063 and
do not invent missing L064 content or assign it extra claims. This is
an edition boundary, not a statement that no other appendix edition
ever existed.

The project contributions are the shared coordinate, transporter and
normalizer arguments; explicit unordered orbit corrections; complete
local stabilizer calculations; and the proofs and certificates identifying
the full graph groups. Where an argument reuses a source proof, such as
the fourth-arc obstruction in the $S_4,D_{24}$ case, that attribution
remains in its component. Distinct review records also disclose the
collaborative contributions to the cubic group bounds and the Frobenius
argument. Exact source ledgers and input bindings accompany the collection.
No compression ratio, new open-conjecture answer or Lean verification is
claimed in this completion. Those tasks remain outside the present intake
stopping boundary.

\begin{thebibliography}{9}
\bibitem{partII} F.~Buekenhout, J.~De Saedeleer and D.~Leemans,
\emph{On the rank two geometries of the groups $\PSL(2,q)$: part II},
Ars Math. Contemp. \textbf{6} (2013), 365--388.
\bibitem{appendix} The same authors, \emph{Part II (Appendix)},
retained 21-page edition.
\bibitem{partI} J.~De Saedeleer and D.~Leemans,
\emph{On the rank two geometries of the groups $\PSL(2,q)$: part I},
Ars Math. Contemp. \textbf{3} (2010), 177--192.
\end{thebibliography}
\end{document}
